\documentclass[10pt,a4paper]{amsart}
\usepackage[margin=19mm]{geometry}
\usepackage{amsmath,amssymb,mathtools}
\usepackage[scr=rsfs,cal=euler]{mathalfa}
\usepackage{xfrac}
\usepackage[unicode,hidelinks,bookmarks=false]{hyperref}
\newcommand{\ie}{i.e.\ }
\newcommand{\cf}{cf.\ }
\newcommand{\resp}{respectively\ }
\newcommand{\Subsection}{\subsection}
\providecommand{\nobreakdash}{\nobreak\hbox{-}}
\setlength{\emergencystretch}{2em}
\multlinegap0pt
\usepackage{bm}
\usepackage{bbm}
\usepackage{dsfont}
\usepackage{xspace}
\usepackage{cancel}
\usepackage{mathtools}
\usepackage{amsthm}
\makeatletter
\@ifundefined{theorem}{\newtheorem{theorem}{Theorem}}{}
\@ifundefined{lemma}{\newtheorem{lemma}[theorem]{Lemma}}{}
\makeatother
\title[The independence and clique cover numbers of the squarefree ]{The independence and clique cover numbers of the squarefree graph}
\author{Boris Alexeev, Dustin G. Mixon, Will Sawin}
\date{Source: arXiv:2507.01928v2 (CC BY 4.0)}
\begin{document}
\maketitle

\noindent\emph{Source and licence.} The independence and clique cover numbers of the squarefree graph. Version arXiv:2507.01928v2 (CC BY 4.0), distributed under the Creative Commons Attribution 4.0 International licence, as recorded in the arXiv licence metadata for this exact version and shown on the abstract page https://arxiv.org/abs/2507.01928v2. Full rights notice: licenses/arxiv-2507.01928-CC-BY-4.0.txt.\par\smallskip

\noindent\textit{Editorial scope.} Counted unit: the Lemma whose source label is \texttt{lem.estimates} in the article (source0.tex lines 513--535, source label \texttt{lem.estimates}), delivered together with the article's own complete original proof of it (source0.tex lines 562--660, one \texttt{proof} environment of 99 lines). No mathematics is written, repaired or re-derived: statement and proof are the article's own text line for line. The proof is detached from the statement in the source: the article states the result at lines 513--535 and prints its proof at lines 562--660, and the proof environment's own header names the statement, so the association is the article's own. Required context retained uncounted: lem.main (lines 367--376); lem.negative moment method (lines 539--547). Every definition, display and label of those lines is retained unchanged. Omitted from this package and not redistributed: 1--366, 377--512, 536--538, 548--561, 661--1034. Nothing inside a retained excerpt is omitted or altered: every retained body is delivered exactly as the article's lines are written, so no omission receipt of any kind accompanies this package. Citations: the retained excerpts carry no citation command at all, so no bibliography entry of the article is transcribed and no reference is redistributed. Reference adaptation: none was needed and none was made; every reference inside the delivered unit resolves inside it, so no unresolved reference or citation is produced. Printed slips kept literal: no printed word, symbol, hypothesis or number of the counted statement or of its proof was altered. Layout and class: the article's own class is \texttt{article} with packages including babel, geometry, amsmath, amssymb, amsthm, graphicx, hyperref, xfrac, tikz, array, pgfplots; the wrapper is the lane's pinned \texttt{amsart} editorial wrapper at 10pt on A4, which restates the theorem-like environments the retained text needs and adds the article's own macro definitions that the retained text needs (none), with extra wrapper packages (bm, bbm, dsfont, xspace, cancel, mathtools). The delivered numbering is the unit's own. Assets: no retained span contains a figure environment, a graphics-inclusion command, a PDF-page inclusion, a table or a TikZ picture; no image file is copied into this package, the package declares no asset dependency and no image file is redistributed with it. Licence: version arXiv:2507.01928v2 of this article is distributed under the Creative Commons Attribution 4.0 International licence, as recorded in the arXiv licence metadata for this exact version and shown on the abstract page https://arxiv.org/abs/2507.01928v2. A full rights notice is delivered at licenses/arxiv-2507.01928-CC-BY-4.0.txt. Characters: every retained excerpt is pure ASCII, so the delivered engine, XeLaTeX with its default Latin Modern text font, reports no missing character.
\begin{lemma}
\label{lem.main}
For every $n\geq10^{10}$, there exists $K=K(n)$ such that the following hold:
\begin{itemize}
\item[(a)]
For each $k\leq K$, there are at most $k$ odd vertices that are each coprime to at most $k$ even vertices.
\item[(b)]
Each odd vertex coprime to $k>K$ even vertices shares a factor with at most $k$ odd vertices.
\end{itemize}
\end{lemma}

\begin{lemma}
\label{lem.negative moment method}
Suppose $X\in[0,b]$ almost surely and $\mathbb{E}[X^{-1}] < \infty$.
Then for every $s\in(0,b)$,
\[
\mathbb{P}\{X\leq s\}
\leq\frac{\mathbb{E}[X^{-1}]-b^{-1}}{s^{-1}-b^{-1}}.
\]
\end{lemma}

\begin{lemma}
\label{lem.estimates}
For every $n\geq 10^{10}$, the following hold:
\begin{itemize}
\item[(a)]
$\#\{\ell\in V(n):2\nmid\ell\}
=\big(1+E(\varepsilon_a)\big)\cdot\frac{2}{3}v(n)$ with $\varepsilon_a:=8.5\cdot10^{-5}$,
\item[(b)]
$\#\{\ell\in V(n):2\nmid\ell,\,3\nmid\ell\}
=\big(1+E(\varepsilon_b)\big)\cdot\frac{1}{2}v(n)$ with $\varepsilon_b:=1.8\cdot10^{-4}$,
\item[(c)]
$F(\ell,n)
=\big(1+E(\varepsilon_c)\big)\cdot f(\ell)$ for each odd $\ell\leq n$ with $\varepsilon_c:=3.8\cdot10^{-3}$,
\item[(d)]
$\sum_{\ell\in V(n),\,2\nmid\ell,\,3\nmid\ell}f(\ell)^{-1}
=\big(1+E(\varepsilon_d)\big)\cdot\frac{36}{91\zeta(3)}\cdot n$ with $\varepsilon_d:=1.3\cdot10^{-3}$,
\item[(e)]
$\sum_{\ell\in V(n),\,2\nmid\ell,\,3\mid\ell}f(\ell)^{-1}
=\big(1+E(\varepsilon_e)\big)\cdot\frac{16}{91\zeta(3)}\cdot n$ with $\varepsilon_e:=2.2\cdot10^{-3}$,
\end{itemize}
where $E(\varepsilon)$ denotes an error term whose absolute value is at most~$\varepsilon$.
Similar to the big~O notation $O(\varepsilon)$, the exact value of $E(\varepsilon)$ is possibly different with each use.
\end{lemma}

\begin{proof}[Proof of Lemma~\ref{lem.main}]
We use Lemma~\ref{lem.estimates} to verify the result with
\[
K
:=0.672\cdot c(1,n).
\]
(Notably, $0.672$ is the notion of $\frac{2}{3}+\varepsilon$ that works well with our nonasymptotic analysis.)

For (a), draw $\ell$ uniformly from the odd vertices.
We wish to show that for every $k\leq K$,
\[
\mathbb{P}\{C(\ell,n)\leq k\}
\leq \frac{k}{\#\{\ell\in V(n):2\nmid\ell\}}.
\]
To accomplish this, we apply Lemma~\ref{lem.negative moment method} after conditioning on whether $3$ divides $\ell$.
(Without conditioning, the resulting bound is too weak for our purposes.)
Conditioning gives
\[
\mathbb{P}\{C(\ell,n)\leq k\}
=\mathbb{P}_{3\mid\ell}\bigg\{F(\ell,n)\leq\frac{k}{c(1,n)}\bigg\}\cdot\mathbb{P}\{\,3\mid\ell\,\}+\mathbb{P}_{3\nmid\ell}\bigg\{F(\ell,n)\leq\frac{k}{c(1,n)}\bigg\}\cdot\mathbb{P}\{\,3\nmid\ell\,\}.
\]
Lemma~\ref{lem.estimates}(a) and (b) together give
\[
\mathbb{P}\{\,3\nmid\ell\,\}
=\frac{\#\{\ell\in V(n):2\nmid\ell,\,3\nmid\ell\}}{\#\{\ell\in V(n):2\nmid\ell\}}
\leq\frac{(1+\varepsilon_b)\cdot\frac{1}{2}v(n)}{(1-\varepsilon_a)\cdot\frac{2}{3}v(n)}
=\frac{3(1+\varepsilon_b)}{4(1-\varepsilon_a)},
\]
\[
\mathbb{P}\{\,3\mid\ell\,\}
=1-\mathbb{P}\{\,3\nmid\ell\,\}
\leq1-\frac{3(1-\varepsilon_b)}{4(1+\varepsilon_a)}.
\]
Considering $f(\ell)\leq\frac{3}{4}$ whenever $2\nmid\ell$ and $3\mid\ell$, Lemma~\ref{lem.estimates}(c) gives that $F(\ell,n)\leq\frac{3}{4}\cdot(1+\varepsilon_c)$ for all such $\ell$.
Then Lemma~\ref{lem.negative moment method} implies
\[
\mathbb{P}\{C(\ell,n)\leq k\}
\leq\frac{\mathbb{E}_{3\mid\ell}[F(\ell,n)^{-1}]-(\frac{3}{4}(1+\varepsilon_c))^{-1}}{(\frac{k}{c(1,n)})^{-1}-(\frac{3}{4}(1+\varepsilon_c))^{-1}}\cdot\bigg(1-\frac{3(1-\varepsilon_b)}{4(1+\varepsilon_a)}\bigg)+\frac{\mathbb{E}_{3\nmid\ell}[F(\ell,n)^{-1}]-1}{(\frac{k}{c(1,n)})^{-1}-1}\cdot\frac{3(1+\varepsilon_b)}{4(1-\varepsilon_a)}.
\]
Next, we apply Lemma~\ref{lem.estimates}(c), and then (b) and (d) to get
\begin{align*}
\mathbb{E}_{3\nmid\ell}[F(\ell,n)^{-1}]
&\leq\frac{1}{1-\varepsilon_c}\cdot\mathbb{E}_{3\nmid\ell}[f(\ell)^{-1}]\\
&=\frac{1}{1-\varepsilon_c}\cdot\frac{1}{\#\{\ell\in V(n):2\nmid\ell,\,3\nmid\ell\}}\sum_{\substack{\ell\in V(n)\\2\nmid\ell,\,3\nmid\ell}}f(\ell)^{-1}
\leq\frac{1+\varepsilon_d}{(1-\varepsilon_c)(1-\varepsilon_b)}\cdot\frac{12\pi^2}{91\zeta(3)}.
\end{align*}
We similarly bound $\mathbb{E}_{3\mid\ell}[F(\ell,n)^{-1}]$ after applying Lemma~\ref{lem.estimates}(a) and (b) to estimate the underlying sample space:
\begin{align*}
\#\{\ell\in V(n):2\nmid\ell,\,3\mid\ell\}
&=\#\{\ell\in V(n):2\nmid\ell\}-\#\{\ell\in V(n):2\nmid\ell,\,3\nmid\ell\}\\
&\geq(1-\varepsilon_a)\cdot\frac{2}{3}v(n)-(1+\varepsilon_b)\cdot\frac{1}{2}v(n)\\
&=(1-4\varepsilon_a-3\varepsilon_b)\cdot\frac{1}{6}v(n).
\end{align*}
Combining this with Lemma~\ref{lem.estimates}(c) and (e) then gives
\[
\mathbb{E}_{3\mid\ell}[F(\ell,n)^{-1}]
\leq\frac{1}{1-\varepsilon_c}\cdot\frac{1}{\#\{\ell\in V(n):2\nmid\ell,\,3\mid\ell\}}\sum_{\substack{\ell\in V(n)\\2\nmid\ell,\,3\mid\ell}}f(\ell)^{-1}
\leq\frac{1+\varepsilon_e}{(1-\varepsilon_c)(1-4\varepsilon_a-3\varepsilon_b)}\cdot\frac{16\pi^2}{91\zeta(3)}.
\]
Denoting $t:=k/c(1,n)=3k/v(n)$, then
\begin{align*}
\mathbb{P}\{C(\ell,n)\leq k\}
&\leq\frac{\frac{1+\varepsilon_e}{(1-\varepsilon_c)(1-4\varepsilon_a-3\varepsilon_b)}\cdot\frac{16\pi^2}{91\zeta(3)}-(\frac{3}{4}(1+\varepsilon_c))^{-1}}{t^{-1}-(\frac{3}{4}(1+\varepsilon_c))^{-1}}\cdot\bigg(1-\frac{3(1-\varepsilon_b)}{4(1+\varepsilon_a)}\bigg)\\
&\qquad+\frac{\frac{1+\varepsilon_d}{(1-\varepsilon_c)(1-\varepsilon_b)}\cdot\frac{12\pi^2}{91\zeta(3)}-1}{t^{-1}-1}\cdot\frac{3(1+\varepsilon_b)}{4(1-\varepsilon_a)}.
\end{align*}
One may verify that for all $t\leq 0.672$, the right-hand side is less than $\frac{t}{2(1+\varepsilon_a)}$.
As such, for every $k\leq K:=0.672\cdot c(1,n)$, it holds that 
\[
\mathbb{P}\{C(\ell,n)\leq k\}
\leq \frac{t}{2(1+\varepsilon_a)}
= \frac{k}{(1+\varepsilon_a)\cdot\frac{2}{3}v(n)}
\leq\frac{k}{\#\{\ell\in V(n):2\nmid\ell\}},
\]
where the last step applies Lemma~\ref{lem.estimates}(a).

For (b), consider any odd vertex $\ell$ such that $C(\ell,n)> K:=0.672\cdot c(1,n)$. 
Note that by doubling the odd vertices coprime to $\ell$, we obtain the even squarefree numbers $\leq 2n$ that are coprime to $\ell$, of which there are $C(\ell,2n)$.
Also note that
\[
c(1,2n)
=\frac{1}{3}v(2n)
=\frac{2}{3}v(n)
=2c(1,n).
\]
It follows that the number of odd vertices that share a factor with $\ell$ is
\[
\begin{aligned}
&\#\{\ell'\in V(n):2\nmid\ell\}-C(\ell,2n)\\
&\qquad=\#\{\ell'\in V(n):2\nmid\ell\} - c(1,2n) F( \ell, 2n )&&\text{($F( \ell, 2n )=\tfrac{C(\ell,2n)}{c(1,2n)}$)} \\
&\qquad\leq(1+\varepsilon_a)\cdot\tfrac{2}{3}v(n) - c(1,2n)\cdot(1-\varepsilon_c)f(\ell)\qquad&&\text{(Lemma~\ref{lem.estimates}(a) and (c))}\\
&\qquad=2\big(1+\varepsilon_a-(1-\varepsilon_c)\cdot f(\ell)\big)\cdot c(1,n)&&\text{($c(1,n)=\tfrac{1}{3}v(n)$, $c(1,2n)=2c(1,n)$)}\\
&\qquad\leq2\big(1+\varepsilon_a-\tfrac{1-\varepsilon_c}{1+\varepsilon_c}\cdot F(\ell,n)\big)\cdot c(1,n)&&\text{(Lemma~\ref{lem.estimates}(c))}\\
&\qquad< 2\big(1+\varepsilon_a-\tfrac{1-\varepsilon_c}{1+\varepsilon_c}\cdot 0.672\big)\cdot c(1,n)&&\text{$(F( \ell, n )=\tfrac{C(\ell,n)}{c(1,n)}>0.672$)}\\
&\qquad\leq 0.672\cdot c(1,n)&&\text{($\varepsilon_a=8.5\cdot10^{-5}$, $\varepsilon_c=3.8\cdot10^{-3}$)}\\
&\qquad<C(\ell,n)&&\text{($C(\ell,n)>0.672\cdot c(1,n)$)},
\end{aligned}
\]
as desired.
\end{proof}

\end{document}
